\chapter{Dataset Description and Initial Risk Analysis}
\label{ch:dataset}

\section{Source and acquisition}

The project uses OpenNeuro dataset \code{ds004504}, local version 1.0.9, released
under CC0. It contains routine eyes-closed resting-state recordings from 88
people: 36 \AD{}, 23 \FTD{}, and 29 \CN{}. Data were acquired with a clinical
Nihon Kohden system at 500 Hz using 19 scalp electrodes positioned according to
the international 10--20 system. The accompanying dataset paper describes the
cohort, recording protocol, BIDS structure, MMSE metadata, and author-supplied
preprocessing \citep{miltiadous2023dataset,openneuro2026dataset}.

The local derivative EEGLAB files are used because the dataset authors recommend
their preprocessed release. Those derivatives include approximately 0.5--45 Hz
band limitation, mastoid referencing, Artifact Subspace Reconstruction (ASR),
ICA, and automatic ICLabel-based component removal. ICLabel is an automated
classifier for EEG independent components \citep{piontonachini2019iclabel}. The
project therefore does not repeat ASR or ICA and does not relabel the derivatives
as raw unprocessed recordings.

\section{Participants and demographics}

Table~\ref{tab:source-demographics} reports the local participant inventory.
Age and sex are used to audit selection, not as predictive features. MMSE is
retained as metadata but excluded from all classifiers, preventing an outcome-
adjacent clinical score from inflating EEG-only performance.

\begin{table}[htbp]
\centering
\begin{tabular}{lrrrrr}
\toprule
\textbf{Group} & \textbf{People} & \textbf{Female} & \textbf{Male} & \textbf{Age range} & \textbf{Age mean}\\
\midrule
\AD{} & 36 & 24 & 12 & 49--79 & 66.39\\
\FTD{} & 23 & 9 & 14 & 44--78 & 63.65\\
\CN{} & 29 & 11 & 18 & 57--78 & 67.90\\
\midrule
Total & 88 & 44 & 44 & 44--79 & --\\
\bottomrule
\end{tabular}
\caption{Source participant distribution. Sex coding follows the supplied metadata.}
\label{tab:source-demographics}
\end{table}

\section{Recording-duration coverage}

The source contains 70,590.3 seconds (19.61 hours) of recording. The duration
imbalance is substantial: \FTD{} contributes both the fewest people and the
least total time. This makes \FTD{} the correct reference margin for a controlled
three-class primary cohort.

\begin{table}[htbp]
\centering
\begin{tabular}{lrrrr}
\toprule
\textbf{Group} & \textbf{People} & \textbf{Seconds} & \textbf{Hours} & \textbf{Complete 4-s windows}\\
\midrule
\AD{} & 36 & 29,403.3 & 8.17 & 7,267\\
\FTD{} & 23 & 16,753.4 & 4.65 & 4,137\\
\CN{} & 29 & 24,433.6 & 6.79 & 6,015\\
\midrule
Total & 88 & 70,590.3 & 19.61 & 17,419\\
\bottomrule
\end{tabular}
\caption{Source duration and complete non-overlapping four-second window inventory.}
\label{tab:duration-source}
\end{table}

Across all participants, the longest recording is \code{sub-010} (\AD{},
1,291.1 s) and the shortest is \code{sub-003} (\AD{}, 307.1 s). Within \FTD{},
the range is \code{sub-070} at 479.1 s to \code{sub-074} at 1,026.9 s; within
\CN{}, it is \code{sub-060} at 751.5 s to \code{sub-040} at 1,017.1 s.
Non-destructive duration-ranked views were created under
\path{preprocessing/01_duration_sorted}; the source files themselves were not
moved or renamed.

\begin{figure}[!htbp]
\centering
\begin{tikzpicture}[node distance=8mm and 9mm,
block/.append style={text width=63mm,minimum height=16mm},
process/.append style={text width=63mm,minimum height=16mm}]
\node[block,text width=116mm] (src) at (0,0) {88 source participants: 17,419 complete four-second windows\\First screen: 17,391 approved and 28 flagged};
\node[block] (p1) at (-3.8,-2.4) {Stage 02 primary\\12,396 windows, 69 people\\4,132 windows per diagnosis};
\node[block] (left) at (3.8,-2.4) {Stage 02 leftover\\5,023 windows retained\\Includes the 28 first-screen flags};
\node[process,below=of p1] (car) {Controlled-primary CAR and QC\\65 signal exclusions; 61 balance trims\\12,270 retained (4,090 per class)};
\node[process,below=of left] (rejoin) {Later full-cohort analysis\\Uniform CAR, QC and boundary guards\\on both Stage 02 partitions};
\node[block,below=of car] (v2) {V2 boundary-safe primary\\11,942 windows, 69 people\\328 boundary-intersecting windows removed};
\node[block,below=of rejoin] (full) {V3 / V4 full cohort\\16,824 windows, 88 people\\Per-person pooling and class weights};
\draw[arrow] (src.south) -- ++(0,-0.3) -| (p1.north);
\draw[arrow] (src.south) -- ++(0,-0.3) -| (left.north);
\draw[arrow] (p1)--(car);
\draw[arrow] (car)--(v2);
\draw[arrow] (left)--(rejoin);
\draw[arrow] (p1.east) -- ++(0.4,0) |- (rejoin.west);
\draw[arrow] (rejoin)--(full);
\end{tikzpicture}
\caption{Dataset lineage. The controlled-primary path and the later full-cohort path create separate analysis indexes. Both Stage 02 partitions are preserved.}
\label{fig:dataset-flow}
\end{figure}


\section{Initial methodological risks}

\subsection{Class count versus window count}

The 17,419 complete windows are not 17,419 independent patients. Windows from
one recording share anatomy, acquisition conditions, diagnosis, and temporal
context. A random window split would permit near-duplicate subject signatures
to appear in both training and test sets. Every valid experiment must therefore
split by subject and pool predictions or features at the subject level.

\subsection{Duration confounding}

If all windows are used naively, \AD{} supplies 1.76 times as many complete
windows as \FTD{}. A classifier may then learn recording exposure or obtain an
overly stable estimate for long recordings. The controlled primary cohort uses
duration buckets and equal windows; the later full cohort averages windows
within person and applies diagnosis-level sample weights.

\subsection{Demographic mismatch}

\FTD{} is younger on average and has a different sex ratio. Exact demographic
identity is neither feasible nor desirable with 23 people per class. The design
uses broad matching for the controlled cohort, records remaining differences,
and excludes age and sex from model predictors.

\subsection{Reference and channel-order dependence}

Graph methods are sensitive to electrode labels and coordinate order. The first
tensorizer used a hard-coded order that did not match the EEGLAB row order.
Later code reads and verifies \code{chanlocs} for every subject. This correction
changes anatomical interpretation, not channel-permutation-invariant operations
such as PCA, global statistics, or common-average reference.

\subsection{Boundary discontinuities}

EEGLAB \code{boundary} events indicate data cuts. Four-second windows near such
events can contain artificial discontinuities even after amplitude QC. V2 and
later analyses therefore remove any window intersecting a fixed 0.5 s guard on
either side of each boundary.

